\documentclass[11pt]{article}
\usepackage[margin=1in]{geometry}
\usepackage{amsmath,amssymb,booktabs}

\newcommand{\E}{\mathbb{E}}
\emergencystretch=3em

\title{Place fields from a rotation of the ideal code}
\author{}
\date{}

\begin{document}
\maketitle

\noindent A rotation of the ideal code changes which functions of position the units carry, not
the similarity between positions. Each rotated unit is the same $M$ waves with its own amplitudes
and phases. Setting every phase so that all waves peak at one position $c$ gives a place field
centred on $c$. With equal amplitudes the field is the kernel itself, with standard deviation $r$,
peak $1$, plus a weak ripple over the whole arena. A full rotation can give every one of the $2M$
units such a field, with centres about $3.35r$ apart and peaks about $0.87$, but each field then
holds only about a fifth of its unit's energy and the ripple holds the rest. A place code without
ripples would need about $7{,}300$ units, seven times the $2M$ the ideal has. Checks are in
\texttt{analysis/hopfield\_probe/kernel\_scaling/place\_fields\_rotation.py}
($L = 1716$, $r = 16$, $M = 512$).

\section{Setup}

The arena is a 2D torus of side $L$ cells. The ideal code has $M$ waves with integer frequencies
$n_k$, drawn from a Gaussian with standard deviation $\sigma_n = L/(2\pi r)$ per axis, and angular
frequencies $\omega_k = 2\pi n_k/L$, so $\omega \sim \mathcal{N}(0, \sigma_\omega^2 I)$ with
$\sigma_\omega = 1/r$. Here the frequencies are distinct up to sign and nonzero, so the $2M$
coordinates are orthogonal functions on the torus. The code and its kernel are
\begin{gather}
  z(p) = \frac{1}{\sqrt M}\bigl(\cos\omega_k\!\cdot p,\ \sin\omega_k\!\cdot p\bigr)_{k=1}^M, \\
  K_M(\Delta) = z(p)\cdot z(p+\Delta) = \frac1M\sum_k \cos(\omega_k\!\cdot\Delta) \approx
  k(\Delta) = e^{-|\Delta|^2/2r^2}.
\end{gather}
It is convenient to define
\begin{equation}
  D = \frac{L^2}{\pi r^2} \approx 3{,}660,
\end{equation}
the effective number of dimensions of the exact Gaussian kernel on this arena, derived in
\texttt{kernel\_error.tex}.

\section{A rotation recombines the waves}

A rotation is $y(p) = Q\,z(p)$ with $Q$ orthogonal. It keeps every similarity,
$y(p)\cdot y(q) = z(p)\cdot z(q) = K_M(p - q)$. Unit $i$ is row $q_i = (a_{ik}, b_{ik})_k$ of $Q$
applied to the code:
\begin{equation}
  y_i(p) = \frac{1}{\sqrt M}\sum_k \bigl(a_{ik}\cos\omega_k\!\cdot p + b_{ik}\sin\omega_k\!\cdot p\bigr)
         = \frac{1}{\sqrt M}\sum_k \rho_{ik}\cos(\omega_k\!\cdot p - \varphi_{ik}),
  \qquad \sum_k \rho_{ik}^2 = 1 .
  \label{eq:unit}
\end{equation}
So every rotated unit is a sum of the same $M$ waves. The row chooses an amplitude $\rho_{ik}$ and a
phase $\varphi_{ik}$ for each wave, and orthogonality of $Q$ couples the choices across rows.

Two identities hold for every rotation. Since $|z(p)| = 1$ and $Q$ is orthogonal,
\begin{equation}
  \sum_i y_i(p)^2 = 1 \ \text{at every } p, \qquad |y_i(p)| \le 1 .
  \label{eq:norm}
\end{equation}
Since distinct waves are orthogonal on the torus and $\sum_p \cos^2 = \sum_p \sin^2 = L^2/2$,
\begin{equation}
  \sum_p y_i(p)^2 = \frac{L^2}{2M}\sum_k\rho_{ik}^2 = \frac{L^2}{2M}\quad\text{for every unit.}
  \label{eq:energy}
\end{equation}
Call this the energy of a unit. It does not depend on $Q$.

\section{One place field: align the phases}

Choose the phases so that every wave peaks at a centre $c$, $\varphi_{ik} = \omega_k\cdot c$. Then
\begin{equation}
  y(p) = \frac{1}{\sqrt M}\sum_k \rho_k\cos\bigl(\omega_k\!\cdot(p - c)\bigr),\qquad
  y(c) = \frac{1}{\sqrt M}\sum_k\rho_k \le 1,
\end{equation}
by Cauchy--Schwarz, with equality only for equal amplitudes $\rho_k = 1/\sqrt M$. Equal amplitudes
and aligned phases make the row equal to the code of the centre, $q = z(c)$, and the unit's tuning
is the kernel centred on $c$:
\begin{equation}
  y_c(p) = z(c)\cdot z(p) = K_M(p - c) = k(p - c) + \varepsilon(p - c).
  \label{eq:matched}
\end{equation}

Why this is a field. At $p = c$ every wave is at its peak. Moving a displacement $\Delta$ away
advances wave $k$ by the phase $\omega_k\cdot\Delta$. Averaged over the frequency distribution,
\begin{equation}
  \frac1M\sum_k\cos(\omega_k\!\cdot\Delta) \approx \E_\omega\cos(\omega\cdot\Delta)
  = e^{-\sigma_\omega^2|\Delta|^2/2},
\end{equation}
so the waves stay in step for $|\Delta| \lesssim 1/\sigma_\omega$ and then drift out of step and
cancel. The field is the Fourier transform of the frequency distribution, and its width is the
distance over which the waves stay in step:
\begin{equation}
  \text{field sd} = \frac{1}{\sigma_\omega} = r .
\end{equation}
Unequal amplitudes reweight the frequency distribution. Weighting high frequencies more makes the
field narrower, down to roughly the shortest wavelength in the code, and weighting low frequencies
more makes it wider. Either way the peak falls below $1$.

The ripple $\varepsilon$ is the part of the cancellation that $M$ waves cannot complete. Its size
follows from the energy \eqref{eq:energy}. The field holds $\sum_p k^2 = \pi r^2$, so
\begin{equation}
  \text{fraction of energy in the field} = \frac{\pi r^2}{L^2/2M} = \frac{2M}{D},\qquad
  \text{ripple rms} = \sqrt{\frac{1 - 2M/D}{2M}},
\end{equation}
valid while $2M \ll D$. At $M = 512$ the field holds $28\%$ of the energy and the ripple has rms
$0.0265$, spread over the whole arena.

\noindent Check (one row $z(c)$): peak $1.000$; half-maximum area equivalent to a standard
deviation of $15.3$ cells against $r = 16$ (drawing distinct frequencies flattens the spectrum
slightly); energy $2875.6 = L^2/2M$, of which $25\%$ lies within $3r$ of the centre (predicted
$28\%$); ripple rms $0.027$ (predicted $0.0265$); largest ripple value $0.11$.

The same budget bounds any unit whose field has the kernel's width. A field with peak $h \le 1$
and standard deviation $r$ holds $\pi h^2 r^2$, so
\begin{equation}
  \text{share of a unit's energy in a field of width } r \;\le\; \frac{2M}{D} = 0.28,
\end{equation}
and at least $72\%$ of every such unit is ripple.

\section{Several place fields in one rotation}

Rows $z(c_1), \dots, z(c_N)$ can all belong to one rotation only if they are orthonormal:
\begin{equation}
  z(c_i)\cdot z(c_j) = K_M(c_i - c_j) = k(c_i - c_j) + \varepsilon(c_i - c_j) = \delta_{ij}.
\end{equation}
This needs two things.
\begin{itemize}
  \item The fields must not overlap, $k(c_i - c_j) \approx 0$, so centres must be at least about
  $3r$ apart ($k(3r) = 0.011$).
  \item The ripples must not add up. Beyond the fields, the overlaps $\varepsilon(c_i - c_j)$ are
  like random numbers with standard deviation $1/\sqrt{2M}$, so the rows behave like $N$ random
  unit vectors in $2M$ dimensions. Their Gram matrix has eigenvalues in
  $\bigl[(1 - \sqrt{N/2M})^2,\ (1 + \sqrt{N/2M})^2\bigr]$ (Marchenko--Pastur), close to $1$ only
  when $N \ll 2M$.
\end{itemize}
The orthonormal rows closest to $z(c_1), \dots, z(c_N)$ are $Q = G^{-1/2}Z$, where $Z$ stacks the
rows and $G = ZZ^\top$. When $N \ll 2M$, $Q \approx Z - \tfrac12(G - I)Z$, so each unit keeps its
field with peak close to $1$ and picks up small bumps of height about $\tfrac12|\varepsilon| \sim
0.015$ at the other centres. The correction grows with $N$, but the fields survive all the way to a
full rotation, $N = 2M$, with one field per unit:
\begin{center}
\small
\setlength{\tabcolsep}{5pt}
\begin{tabular}{rrcccccc}
  \toprule
  $N$ & spacing & eig.\ of $G$ & MP range & peak & field sd & in field & beyond $3r$ \\
  \midrule
  16 & $26.8r$ & 0.80--1.25 & 0.77--1.27 & 0.997 & 15.3 & 0.25 & 0.11 \\
  64 & $13.4r$ & 0.56--1.44 & 0.56--1.56 & 0.993 & 15.4 & 0.25 & 0.12 \\
  256 & $6.7r$ & 0.00--2.25 & 0.25--2.25 & 0.971 & 15.4 & 0.24 & 0.12 \\
  1024 & $3.35r$ & $2\times10^{-7}$--3.27 & 0--4 & 0.87 & 15.5 & 0.20 & 0.13 \\
  \bottomrule
\end{tabular}
\end{center}
MP is the Marchenko--Pastur range, ``in field'' the share of a unit's energy within $3r$ of its
centre, and ``beyond $3r$'' the largest $|y|$ farther out. The first three rows use centres on an
exact lattice. That makes $G$ a circulant whose eigenvalues
count the frequencies in each residue class, so a class with no frequency gives a zero eigenvalue
at $N = 256$, and at $N = 1024$ the exact lattice leaves $G$ singular. The last row therefore
jitters each centre by up to a quarter of the spacing. $Q$ is then orthogonal to $10^{-12}$, and
over 96 units the peak ranges from $0.86$ to $0.88$ (5th to 95th percentile) and the largest value
outside $3r$ stays below $0.15$. Jittering by half the spacing gives the same numbers within
$0.01$, except for a 5th-percentile peak of $0.82$.

So a full rotation turns the units into place cells (every one of the 96 checked): $2M = 1024$
fields, one per unit, with centres $L/\sqrt{2M} \approx 3.35r$ apart, standard deviation about $r$
and peak about $0.87$.
Each field holds about a fifth of its unit's energy, close to the bound $2M/D \times 0.87^2 = 0.21$,
and the rest is a ripple of rms $0.028$ over the whole arena.

\section{What the rotated code is and is not}

Summing \eqref{eq:norm} over the arena gives $\sum_i\sum_p y_i(p)^2 = L^2$ for any code of unit
norm with any number of units. A unit that is a single bump with peak $h \le 1$ and standard
deviation $\sigma$ has energy $\pi h^2\sigma^2 \le \pi\sigma^2$. Two consequences follow.

\paragraph{A place code with this kernel needs many units.} In a place code the similarity between
two positions is the overlap of their field patterns. For Gaussian fields of standard deviation
$\sigma_f$ this overlap is a Gaussian of standard deviation $\sqrt2\,\sigma_f$, so matching the
kernel needs $\sigma_f = r/\sqrt2$. Equivalently, the frequency density $p(\omega)$ of the ideal is
the power spectrum $|\hat f(\omega)|^2$ of the field. Each such field holds at most
$\pi\sigma_f^2 = \pi r^2/2$, so the number of units is at least
\begin{equation}
  N_{\text{place}} \ge \frac{L^2}{\pi\sigma_f^2} = \frac{2L^2}{\pi r^2} = 2D \approx 7{,}300 .
\end{equation}
The ideal has $2M = 1024$ units, and a rotation keeps that number. So no rotation of the ideal at
$M = 512$ is a place code with its kernel.

\paragraph{If every unit were a single bump, the bumps would be wide.} With $2M$ units each a single
bump, $\sum_i \pi h_i^2\sigma_i^2 = L^2$ gives
\begin{equation}
  \sqrt{\langle\sigma^2\rangle} \ge \frac{L}{\sqrt{2\pi M}} \approx 30\ \text{cells},
\end{equation}
set by $L$ and $M$, not by $r$. A tiling of bumps that wide would have population similarity with
standard deviation $\sqrt2 \times 30 \approx 43$ cells, but a rotation keeps it at $r = 16$. So no
rotation makes every unit a single clean bump.

\paragraph{What it is.} The full rotation of Section~4 sits between these two limits. Its fields
have the kernel's width, so each unit's field holds at most $2M/D$ of its energy and the ripple
holds the rest. Averaged over positions, the fields make up about a fifth of each code vector's
squared length. At half maximum they cover about a third of the arena; a position between fields
is coded by the ripples. The fields alone give a similarity of about
$0.2\,e^{-|\Delta|^2/4r^2}$ (the overlap of fields of width $r$), and the ripples supply the rest of
$K_M$. The ripple is what fitting the kernel into $2M < 2D$ dimensions costs.

\section{Summary}

\begin{center}
\begin{tabular}{ll}
  \toprule
  quantity & value \\
  \midrule
  field from one row $z(c)$ & $K_M(p - c)$: sd $r = 1/\sigma_\omega$, peak $1$ \\
  share of energy in a field of width $r$ & $\le 2M/D$ ($0.28$ at $M = 512$) \\
  ripple rms of one row $z(c)$ & $\sqrt{(1 - 2M/D)/2M}$ ($0.027$) \\
  full rotation into place cells & one field per unit, centres $3.35r$ apart, \\
  & sd $\approx r$, peak $\approx 0.87$, field $\approx 20\%$ of energy \\
  place-code field with the same kernel & sd $r/\sqrt2$, power spectrum $= p(\omega)$ \\
  units such a place code needs & $\ge 2L^2/\pi r^2 \approx 7{,}300$ \\
  width if all $2M$ units were single bumps & $\ge L/\sqrt{2\pi M} \approx 30$ cells \\
  \bottomrule
\end{tabular}
\end{center}

\end{document}
